\section{Жадные алгоритмы}

\subsection{Домашнее задание}
\begin{enumerate}
  \item
    \begin{enumerate}
      \item
        Пусть в тексте встречается символ с частотой большей $\frac{2}{5}$.
        Докажите, что \textbf{любой} двоичный беспрефиксный код, который минимизирует длину закодированного текста,
        будет содержать слово длины $1$. \textit{(Любой --- значит не только тот, который получается
        алгоритмом Хаффмана.)}

      \item
        Пусть в тексте частоты всех символов меньше $\frac{1}{3}$.
        Докажите, что \textbf{любой} двоичный беспрефиксный код, который минимизирует длину закодированного текста,
        содержит только кодовые слова длины не менее $2$.
    \end{enumerate}

  \item
 	Преподаватели сделали $n$ заявок на занятие. Каждое
 	занятие начинается в момент $b_i$ и кончается в момент $e_i$
 	(занимает интервал $[b_i, e_i)$). Два занятия в одной аудитории
 	быть не могут. Распределите заявки по аудиториям так, чтобы
 	общее число аудиторий было минимально. Решить за $\O(n \log n)$.

  \item
    $n$ aтлетов хотят выстроить из своих тел башню максимальной
    высоты. Башня это цепочка атлетов, первый стоит на земле, второй
    стоит у него на плечах, третий стоит на плечах у второго и
    т.д. Каждый атлет характеризуется силой $s_i$ и массой $m_i$. Сила
    это максимальная масса, которую атлет способен держать у себя на
    плечах. Известно, что если атлет тяжелее, то он и сильнее, но
    атлеты равной массы могут иметь различную силу. Определите
    максимальную высоту башни. $\O(n \log n)$.

  \item
	Есть $n$ человек. Человек $i$ готов примкнуть к нашей банде, если наш авторитет
	хотя бы $a_i$, при этом он изменит наш авторитет на $b_i$. Наш изначальный авторитет
	равен $A$. $a_i, b_i, A \in \mathbb{Z}$
	\begin{enumerate}
		\item Можем ли завербовать всех людей? $\O(n \log n)$.
		\item Какое максимальное число людей мы можем завербовать? $\O(n \log n)$.
	\end{enumerate}

\subsection*{Дополнительные задачи}

  \item
    Будем называть \textit{независимым множеством} (или \textit{антикликой}) подмножество вершин
    графа, в котором нет ребер. Пусть в графе $G$ есть $V$ вершин и $E$ ребер, а максимальная
    степень равна $d$. Найдите в нём независимое множество размера хотя бы:
	\begin{enumerate}
	  \item $\frac{n}{d+1}$ за время $\O(V + E)$
	  \item $\sum\limits_{v\in V(G)} \frac{1}{\deg(v)+1}$ за время $\O(V \log{V} + E)$
	\end{enumerate}
	Считайте, что граф уже дан в памяти в виде массива, где для каждой вершины хранится список её соседей.

  \item {\bf Пятачок, у тебя есть дома ружье?}\\
	Даны $n$ непересекающихся кругов на плоскости. Мы стоим в точке $(0,0)$ и можем стрелять по прямой.
	Минимальным числом выстрелов проткнуть все круги.

  \item
	$n$ школьников упали в яму глубины $S$. Каждый школьник имеет рост (от ног до плеч) $h_i$
	и длину рук $l_i$. Школьники могут вставать друг другу на плечи, верхний школьник может
	вытянуть руки. Чтобы выбраться из ямы необходимо дотянуться руками до уровня земли.
	\begin{enumerate}
		\item Могут ли выбраться все школьники? $\O(n \log n)$.
		\item Какое максимальное число школьников может выбраться? $\O(n \log n)$.
	\end{enumerate}

  \item
    В фирму поступают заказы, которые можно выполнять в произвольном
    порядке. В каждый момент времени можно работать ровно над одним
    заказом.  Изначально заказов нет, $i$-й заказ поступает в момент
    времени $r_i$, работать над ним нужно $t_i$. Пусть $e_i$~-- момент
    окончания выполнения заказа номер $i$. Распределите работу над
    заказами так, чтобы минимизировать $\sum_i e_i$. Переходить от
    одного заказа к другому можно в любой момент времени (даже если
    заказ не доделан до конца, незаконченный заказ можно будет
    возобновить с того же места).
    Свойства заказа ($r_i, t_i$) не известны до
    момента его поступления. Время $\O(n \log n)$, при условии, что
    всего поступит $n$ заказов.

  \item
    \textit{Раскраской вершин} графа $G = (V, E)$ называется функция $c: V \to [m]$, сопоставляющая каждой вершине $G$ цвет от $1$ до $m$.
 	Раскраска называется \textit{правильной}, если каждая пара соседних вершин имеет разные цвета.

 	Для неориентированного графа $G = (V, E)$ его \textit{хроматическим числом} $\chi(G)$ называется наименьшее возможное число цветов в правильной раскраске $G$.

 	Для графа $G$ обозначим размер его максимального полного подграфа через $\omega(G)$.

 	Рассмотрим на вещественной прямой замкнутые отрезки $I_1, I_2, \dots, I_n$. Сопоставим каждому отрезку $I_i$ вершину $v_i$
 	и каждой паре пересекающихся отрезков $(I_i, I_j)$ ребро $(v_i, v_j)$. Такой граф будем называть \textit{интервальным графом}.

 	Будем называть граф $G$ \textit{совершенным}, если для любого его индуцированного подграфа $H$ верно $\omega(H) = \chi(H)$.

 	Докажите, что каждый интервальный граф совершенен. Приведите алгоритм, красящий интервальный граф $G = (V, E)$ с $|V| = n$ в $\omega(H)$ цветов за время $O(n \log(n))$.

  \item
 	В фирму поступают заказы, которые можно выполнять в произвольном
 	порядке. В каждый момент времени можно работать ровно над одним
 	заказом.  Изначально заказов нет, $i$-й заказ поступает в момент
 	времени $r_i$, работать над ним нужно $t_i$.

 	Все заказы объединены в проекты (один заказ относится к одному
 	проекту, заказы из одного проекта могут поступать не подряд).

 	Пусть $e_i$~-- момент окончания выполнения последнего (в порядке
 	выполнения) из заказов в проекте с номером $i$.

 	Нужно распределить работу над заказами так, чтобы минимизировать
 	$\sum_i e_i$. Свойства заказа ($r_i, t_i$, его проект) не известны
 	до момента его поступления.  Переходить от одного заказа к другому
 	можно в любой момент времени (даже если заказ не доделан до конца).

 	Придумайте решение, которое не более чем в два раза хуже
 	оптимального.  Время $\O(n \log n)$, при условии, что всего поступит $n$ заказов.

\end{enumerate}

\newpage

\subsection{Решение}

\begin{enumerate}

\item \dots

\item

    Создадим модель события, которое привязано к какому-то времени $t_i$ и имеет одно из двух состояний
    $s_i \in \{+1, -1\}$: начало или конец заявки соответственно.

    Каждую заявку можно представить в виде двух событий: \( (b_i, +1), (e_i, -1) \).

    Тогда, множество всех событий задается как:
    \[
        E = \{ (b_1, +1), (e_1, -1), (b_2, +1), (e_2, -1), \dots, (b_n, +1), (e_n, -1) \}
    \]

    Отсортируем множество $E$ по времени:
    \[
       (t_1, s_1) < (t_2, s_2) \iff t_1 < t_2 \lor (t_1 = t_2 \land s_1 < s_2), \text{ где } s_1, s_2 \in \{+1, -1\}
    \]

    Далее подсчитаем максимум активных событий в текущий момент времени.
    Для этого будем итерироваться по событиям и поддерживать текущее количество активных
    событий - инкрементируем счетчик при встрече события начала и декрементируем при
    встрече события конца.

    Максимум из всех значений счетчика и будет ответом.

\item

    Будем жадно составлять башню из атлетов сверху вниз, проверяя, что добавление
    нового атлета не сломает башню - то есть, что новый атлет может удержать всех
    выше стоящих.

    Определим множество атлетов $A = {a_1, a_2, \dots, a_n}$, где $a_i = (w_i, s_i)$ - масса и сила
    $i$-го атлета соответственно.

    Атлеты сортируются по правилу:
    \[
        a_i < a_j \iff (w_i < w_j) \lor (w_i = w_j \land s_i < s_j).
    \]

    Пусть \( T \) — башня из атлетов, \( W_T = \sum_{t \in T} w_t \) — её суммарный вес.
    Для каждого нового атлета \( a_i \) проверяем:
    \[
        s_i \geq W_T.
    \]

    Если условие выполняется:
    \[
        T \gets T \cup \{a_i\}, \quad W_T \gets W_T + w_i, \quad h \gets h + 1.
    \]

    Тем самым за \( O(n) \) мы находим максимальную башню, однако сортировка
    атлетов по массе и силе занимает \( O(n \log n) \).


    %\item
	%Есть $n$ человек. Человек $i$ готов примкнуть к нашей банде, если наш авторитет
	%хотя бы $a_i$, при этом он изменит наш авторитет на $b_i$. Наш изначальный авторитет
	%равен $A$. $a_i, b_i, A \in \mathbb{Z}$
	%\begin{enumerate}
		%\item Можем ли завербовать всех людей? $\O(n \log n)$.
		%\item Какое максимальное число людей мы можем завербовать? $\O(n \log n)$.
	%\end{enumerate}

\item

    Разобъем всех людей на две группы: тех у кого $b_i \geq 0$ и тех у кого $b_i < 0$.
    Это позволит нам рассматривать их отдельно и придти к корректному решению.

    \begin{enumerate}

    \item

        Нетрудно заметить, что если $b_i \geq 0$, то мы можем взять всех таких людей, если они удовлетворяют
        критериям по авторитету.

        Для этого отсортируем всех людей по возрастанию $a_i$ и будем добавлять их в нашу банду.
        Если как минимум один из такой группы не удовлетворяет критериям, то уже на
        этом этапе мы не может завербовать всех людей.

        Далее рассмотрим группу людей с $b_i < 0$.
        Их следует отсортировать по убыванию $b_i$, чтобы уменьшение суммарного авторитета происходило как
        можно медленнее и мы смогли рассмотреть как можно больше людей.
        Если как минимум один из такой группы не удовлетворяет критериям, то мы опять не можем
        взять всех людей.


    \item

        Для того, чтобы максимизировать количество людей, которых мы можем завербовать точно также
        работаем с двумя группами людей.

        Для первой группы выполняем все аналогично - берем всех, которые удовлетворяют критериям.

        Для второй группы заведем приоритетную очередь в которой будем размещать людей на основании
        $a_i + b_i$. Таким образом, мы будем откидывать худших кандидатов, которые уменьшают наш авторитет сильнее всего.
        И также будем обрабатывать их в порядке убывания $b_i$.

        И при вербовке нового человека, если он не удовлетворяет критериям, то мы можем попыться заменить
        худшего кандидата из приоритетной очереди на него, тем самым увеличив общий авторитет.

        За счет этого решения можно покрыть и первый пункт задачи, но это требует дополнительной работы на поддержание
        приоритетной очереди.

    \end{enumerate}

\end{enumerate}
